\documentclass[11pt]{article}
\usepackage[margin=0.95in]{geometry}
\usepackage{amsmath,amssymb,amsthm}
\usepackage[T1]{fontenc}
\usepackage{lmodern}
\usepackage[hidelinks]{hyperref}
\usepackage{microtype}
\setlength{\parindent}{0pt}
\setlength{\parskip}{6pt}
\newtheorem{proposition}{Proposition}
\title{A two-dimensional counterexample to Conjecture 00000009761}
\author{gaochengzhecpu}
\date{October 4, 2026}
\begin{document}
\maketitle
\begin{abstract}
The proposed arithmetic strengthening of the singular-value bound is false,
already for a trace-class operator on $\mathbb C^2$. The upper triangular
matrix $A=\left(\begin{smallmatrix}1&3/2\\0&1\end{smallmatrix}\right)$
has eigenvalues $1,1$ counted with algebraic multiplicity and singular
values $2,1/2$. At $n=2$ the claimed upper bound is $3/4$, whereas
$|\lambda_2|=1$. We give an exact singular value decomposition, a
Gram-matrix calculation, and a Lean 4 verification of the finite-dimensional
counterexample.
\end{abstract}

\section{Statement and conventions}
The conjecture proposes the universal bound
\begin{equation}\label{eq:claim}
 |\lambda_n|\ \leq\ \frac{\bigl(\prod_{k=1}^{n}s_k\bigr)^{1/n}+s_n}{2},
\end{equation}
where the eigenvalues are listed in nonincreasing modulus, with algebraic
multiplicity, and the singular values satisfy $s_1\geq s_2\geq\cdots\geq0$.
The supplied English and Chinese statements contain no normality or
self-adjointness hypothesis. The reference to second-order diagonal models
concerns claimed attainment of optimality, not a restriction on the class
of operators. We disprove the inequality itself, so its asserted universal
optimality and uniqueness cannot hold as stated.

\begin{proposition}
The trace-class operator on the standard complex Euclidean space $\mathbb C^2$
with matrix
\[
 A=\begin{pmatrix}1&3/2\\0&1\end{pmatrix}
\]
violates \eqref{eq:claim} at $n=2$.
\end{proposition}
\begin{proof}
Its characteristic polynomial is $(t-1)^2$, so $\lambda_1=\lambda_2=1$.
The adjoint is the conjugate transpose, and
\[
 A^*A=\begin{pmatrix}1&3/2\\3/2&13/4\end{pmatrix},\qquad
 \det(tI-A^*A)=t^2-\frac{17}{4}t+1
              =(t-4)(t-1/4).
\]
Thus the ordered singular values, the nonnegative square roots of the
eigenvalues of $A^*A$, are $s_1=2$ and $s_2=1/2$.
All finite-dimensional operators are trace class; concretely, the sequence
of singular values padded by zeros is $(2,1/2,0,\ldots)$ and has sum $5/2$.
At $n=2$, the asserted right-hand side becomes
\[
 \frac{\sqrt{2\cdot(1/2)}+1/2}{2}=\frac34<1=|\lambda_2|,
\]
which contradicts \eqref{eq:claim}.
\end{proof}

\newpage
\section{An explicit singular value decomposition}
To independently certify that the listed values are singular values of the
actual matrix, put $c=1/\sqrt5$ and define
\[
 U=c\begin{pmatrix}2&-1\\1&2\end{pmatrix},\qquad
 V=c\begin{pmatrix}1&-2\\2&1\end{pmatrix},\qquad
 S=\begin{pmatrix}2&0\\0&1/2\end{pmatrix}.
\]
Direct multiplication gives $UU^*=VV^*=I$ and
\[
 USV^*=\frac15
 \begin{pmatrix}2&-1\\1&2\end{pmatrix}
 \begin{pmatrix}2&0\\0&1/2\end{pmatrix}
 \begin{pmatrix}1&2\\-2&1\end{pmatrix}
 =\begin{pmatrix}1&3/2\\0&1\end{pmatrix}=A.
\]
Both diagonal entries of $S$ are nonnegative and already in decreasing
order. Hence this is a complete unitary SVD, with no numerical approximation.
The ordinary product bound is not contradicted: for this example
$|\lambda_1\lambda_2|=s_1s_2=1$. The additional averaging with the smaller
singular value is the step that fails. The example is not a normal matrix;
no assertion about a version restricted to normal operators is made here.

\section{Formal verification and its scope}
The project uses Lean 4.19.0 and a pinned Mathlib revision. The matrix in
\texttt{lean/Main.lean} is an actual complex $2\times2$ Mathlib matrix.
The kernel-checked statements establish:
\begin{itemize}
\item the actual characteristic polynomials and full root multisets of
$A$ and $A^*A$, including the repeated eigenvalue;
\item membership of $U,V$ in the actual complex unitary group, the exact
identity $A=USV^*$, and the order and sign of the singular values;
\item the standard Euclidean-space operator and the Mathlib identity
identifying conjugate transpose with the Hilbert-space adjoint;
\item summability of the verified singular values extended by zero, with
exact sum $5/2$, and failure of the proposed $n=2$ inequality.
\end{itemize}
\texttt{OrderedEigenvalues} is defined by the complete characteristic-root
multiset and modulus order. \texttt{OrderedSingularValues} is defined by the
existence of actual unitary matrices giving an SVD with a nonnegative,
ordered diagonal. These are the standard finite-dimensional spectral data.
The final theorem negates the universal inequality for these spectral data.
The theorem
\texttt{full\_counterexample} collects the verified witness and its summable
singular sequence. No general-purpose trace-class API or proof of Lidskii's
theorem is claimed: the finite-dimensional trace-class bridge is precisely
the summability criterion just stated. The $n=2$ root is represented by the
positive real square root.

The supplementary \texttt{verify.py} checks the matrix products,
characteristic polynomials, and strict rational inequality using exact
fractions. Fresh-build logs, a theorem-axiom audit, and the source hashes are
provided in \texttt{verification/}. This proof uses no numerical search.

\small
\textbf{Source.} The Last Math Competition, conjecture 00000009761,
\href{https://github.com/The-Last-Math-Competition/The-Last-Math-Competition/blob/main/conjectures/00000009761.md}{repository statement}.

The submitted \texttt{SOURCE.md} preserves the bilingual snapshot byte for
byte; its source identity is recorded in \texttt{verification/}.
\end{document}
